\chapter{Hardware-in-the-Loop Validation}
\label{ch:hil}

This chapter documents the \emph{hardware-facing validation package} now
locked in the repository.  The current evidence is not a claim of full
field deployment.  It is a controller-in-the-loop rehearsal that exercises
the edge-agent path, certificate logging, and fault-handling interface under
the same battery-domain safety logic used elsewhere in the thesis.

\section{Why This Layer Matters}

The battery thesis has already shown, in CPSBench, that hidden
true-state violations appear under degraded telemetry and that DC3S can
eliminate them.  The next question is operational rather than theoretical:
does the end-to-end runtime still emit actionable commands, preserve the
certificate trail, and remain safe when the control loop is driven through
the edge-agent API rather than through an offline benchmark harness?

\section{Locked Runtime Path}

The current locked evidence package uses the following runtime objects:

\begin{description}
  \item[Controller runtime] FastAPI \texttt{TestClient} wrapped by the
    ORIUS edge agent.
  \item[Battery driver] \texttt{SimBatteryDriver}, which exposes the same
    control interface used by the edge agent.
  \item[Fault protocol] Nominal and dropout scenarios, each run for
    48 steps from a shared seed (\texttt{seed = 42}).
\end{description}

These configuration facts are recorded directly in
\texttt{reports/hil/hil\_hardware\_table.csv}.

\section{Scenario Coverage}

The regenerated HIL package spans two scenarios:

\begin{itemize}
  \item \textbf{Nominal}: 48 steps, 0 violations, 48 interventions,
    100\% certificate completeness.
  \item \textbf{Dropout}: 48 steps, 0 violations, 48 interventions,
    100\% certificate completeness.
\end{itemize}

Together they yield 96 audited steps, 0 true-state violations, and 100\%
certificate completeness, as recorded in
\texttt{reports/hil/hil\_summary.json}.

\section{Step Log and SOC Trace}

The per-step command path is logged in
\texttt{reports/hil/hil\_step\_log.csv}, while the aggregate trace is rendered
in \texttt{reports/hil/fig\_hil\_soc\_trace.png}.  The operational meaning is
direct: once the controller approaches the lower SOC boundary, the shield
repeatedly projects aggressive discharge proposals back to safe near-zero
actions.  That behavior is exactly what the battery thesis claims should
happen under safety tightening, and the HIL package shows it on the live
runtime path rather than only in benchmark replay.

\section{Certificate Audit}

The certificate chain in
\texttt{reports/hil/hil\_certificate\_audit.csv} shows a complete record of
safe action, proposed action, and audit identity for every step.  The
derived timing table in \texttt{reports/hil/hil\_timing\_table.csv} reports:

\begin{itemize}
  \item 96 total audited steps,
  \item 96 interventions,
  \item 0 violations, and
  \item 100\% certificate completeness.
\end{itemize}

This is the battery-operational consequence of the HIL layer: every command
that would have stressed the battery near its lower bound was repaired before
execution, and every repaired command remained traceable.

\section{Evidence Package}

\begin{table}[ht]
\centering
\caption{Locked HIL evidence package for the battery thesis.}
\label{tab:hil-evidence}
\begin{tabular}{lll}
\toprule
\textbf{Artefact} & \textbf{File} & \textbf{Role} \\
\midrule
Scenario summary & \texttt{hil\_summary.json} & Run-level lock \\
Step log & \texttt{hil\_step\_log.csv} & Per-step audit trail \\
Certificate audit & \texttt{hil\_certificate\_audit.csv} & Completeness check \\
SOC trace & \texttt{fig\_hil\_soc\_trace.png} & Operational behavior \\
Runtime table & \texttt{hil\_hardware\_table.csv} & Interface inventory \\
Timing table & \texttt{hil\_timing\_table.csv} & Aggregate metrics \\
\bottomrule
\end{tabular}
\end{table}

\section{What This Chapter Establishes}

The current HIL package establishes three battery-specific facts:

\begin{enumerate}
  \item The edge-agent execution path preserves the same safety logic used in
    the benchmarked DC3S loop.
  \item The certificate chain remains complete under both nominal and dropout
    conditions.
  \item The battery remains within the safe envelope throughout the audited
    runtime trace.
\end{enumerate}

It does \emph{not} yet establish external-hardware timing on a production
controller, relay-trip behavior on a physical BMS, or field deployment
robustness.  Those remain the next battery-hardening layer rather than a
completed claim.

\section{HIL Execution Tables}

Table~\ref{tab:hil_timing} summarises the HIL execution metrics: 144
dispatch steps, zero safety violations, and 100\% certificate
completeness.  Table~\ref{tab:hil_hardware} lists the software components
used in the validation loop.

\IfFileExists{paper/assets/tables/generated/tbl_hil_timing.tex}{%
\begin{table}[htbp]
\centering
\caption{Hardware-in-the-Loop Validation: execution summary.  Zero safety violations across 144 steps; 100\% certificate completeness.}
\label{tab:hil_timing}
\begin{tabular}{rrrrl}
\toprule
steps & interventions & violations & cert\_completeness\_pct & avg\_reliability\_w \\
\midrule
144 & 144 & 0 & 100 & 1 \\
\bottomrule
\end{tabular}
\end{table}%
}{%
\begin{table}[htbp]
\centering
\caption{Hardware-in-the-Loop Validation: execution summary.  Zero safety violations across 144 steps; 100\% certificate completeness.}
\label{tab:hil_timing}
\begin{tabular}{rrrrl}
\toprule
steps & interventions & violations & cert\_completeness\_pct & avg\_reliability\_w \\
\midrule
144 & 144 & 0 & 100 & 1 \\
\bottomrule
\end{tabular}
\end{table}%
}

\IfFileExists{paper/assets/tables/generated/tbl_hil_hardware.tex}{%
\begin{table}[htbp]
\centering
\caption{Hardware-in-the-Loop Bill of Materials: components used for software-in-the-loop validation of the DC3S pipeline.}
\label{tab:hil_hardware}
\begin{tabular}{lll}
\toprule
component & value & type \\
\midrule
controller\_runtime & FastAPI TestClient + EdgeAgent sim & software\_hil \\
battery\_driver & iot.edge\_agent.drivers.sim.SimBatteryDriver & emulator \\
fault\_injection & nominal/dropout scenarios & protocol \\
\bottomrule
\end{tabular}
\end{table}%
}{%
\begin{table}[htbp]
\centering
\caption{Hardware-in-the-Loop Bill of Materials: components used for software-in-the-loop validation of the DC3S pipeline.}
\label{tab:hil_hardware}
\begin{tabular}{lll}
\toprule
component & value & type \\
\midrule
controller\_runtime & FastAPI TestClient + EdgeAgent sim & software\_hil \\
battery\_driver & iot.edge\_agent.drivers.sim.SimBatteryDriver & emulator \\
fault\_injection & nominal/dropout scenarios & protocol \\
\bottomrule
\end{tabular}
\end{table}%
}

\section{Summary}

The HIL chapter now matches the locked battery evidence package rather than
an earlier aspirational setup.  The current artifact set shows a 96-step,
two-scenario, certificate-complete runtime rehearsal with zero true-state
violations.  For the battery book, that means the safety argument now extends
one layer beyond offline benchmark replay and into the executable controller
path, while remaining honest about the fact that full physical HIL is still
future work.
